\par\medskip\noindent
% Continuación Canon: SGA-SRC-001252.
1.1.3.  La estructura de bitorsor de 1.1.2 bajo
$(\sisi{G_P}{G_P},\sisi{G_P}{G_P})$ sobre
$E$, en términos de la estructura de extensión (1.1.1), solo utilizaba
una parte de la estructura multiplicativa de $E$, a saber, la
multiplicación de pares de «puntos» uno de los cuales, al menos,
provenía de un punto de $G$ (por medio de $i$). Vamos a interpretar
ahora la estructura multiplicativa de $E$ como una estructura suplemen
%%%%%%%%%%  scan idx 148  |  folio 137  %%%%%%%%%%
taria sobre el bitorsor considerado. Para ello, consideremos primero
dos secciones $p$, $p'$ de $P$, de donde, con las notaciones de
1.1.2.1, se obtienen dos bitorsors $E_p$, $E_{p'}$ bajo $(G,G)$, a
saber, las imágenes inversas en $E$ de dichas secciones
$p,p' : e \longrightarrow P$ (NB: $e$ denota el objeto final de
$\underline{T}$). Es entonces trivial que la multiplicación de $E$
induce un morfismo de objetos de $\underline{T}$

\begin{equation*}
\varphi : \quad \sisi{P_p}{E_p} \times \sisi{P_{p'}}{E_{p'}}
\ \longrightarrow\ \sisi{P_{pp'}}{E_{pp'}}
\end{equation*}
que satisface, en términos de las estructuras izquierdas y derechas de
los tres bitorsors en juego (y tomando, como de costumbre, puntos con
valores en un objeto arbitrario
$S \in \mathrm{ob}\ \underline{T}$),

\begin{equation}\tag{1.1.3.0}
\left\{\ \begin{array}{l} \varphi(gx,x') = g\ \varphi(x,x')\ ,\ \ \varphi(x,x'g') = \varphi(x,x')g' \\[3ex] \qquad\qquad\qquad \varphi(xg,x') = \varphi(x,gx') \end{array}\right.
\end{equation}
$(g,g' \in G(S)$, $x \in \sisi{P_p}{E_p}(S)$,
$x' \in \sisi{P_{p'}}{E_{p'}}(S))$. La tercera de estas relaciones se
interpreta, en términos del producto fibrado contraído [2],

\begin{equation*}
\sisi{P_p}{E_p} \overset{G}{\times} \sisi{P_{p'}}{E_{p'}} \qquad ,
\end{equation*}
diciendo que el morfismo $\varphi$ se factoriza como

\begin{equation*}
\sisi{P_p}{E_p} \times \sisi{P_{p'}}{E_{p'}} \ \longrightarrow\
\ \sisi{P_p}{E_p} \overset{G}{\times} \sisi{P_{p'}}{E_{p'}}
\ \xrightarrow{\ \varphi_{p,p'}\ }\ \sisi{P_{pp'}}{E_{pp'}}
\qquad ,
\end{equation*}
donde la primera flecha es el morfismo canónico de paso al cociente.
Por otra parte, las dos primeras relaciones significan que el morfismo

\begin{equation}\tag{1.1.3.1}
\varphi_{p,p'} \ :\ \sisi{P_p}{E_p} \overset{G}{\times}
\sisi{P_{p'}}{E_{p'}} \ \longrightarrow\ \sisi{P_{pp'}}{E_{pp'}}
\end{equation}
que acabamos de obtener es un homomorfismo de \underline{bitorsors},
teniendo en cuenta la estructura de bitorsor sobre el primer miembro
para la cual las operaciones
%%%%%%%%%%  scan idx 149  |  folio 138  %%%%%%%%%%
izquierdas (resp. derechas) de $G$ provienen de las operaciones
izquierdas (resp. derechas) de $G$ sobre el primer factor
$\sisi{P_p}{E_p}$ (resp. sobre el segundo factor
$\sisi{P_{p'}}{E_{p'}}$). Observemos que (como todo homomorfismo de
bitorsors bajo un par de grupos dado) el homomorfismo
$\varphi_{p,p'}$ es necesariamente un \underline{isomorfismo}.

La definición del isomorfismo precedente (1.1.3.1) se extiende de
inmediato al caso en que $p,p'$, en lugar de ser secciones de $P$, son
puntos de $P$ con valores en un objeto cualquiera $S$ de
$\underline{T}$: basta aplicar lo anterior al topos inducido
$\underline{T}_{/S}$ y a la extensión de grupos sobre $S$ deducida de
(1.1.1) por cambio de base. Se comprueba entonces que la formación de
los isomorfismos $\varphi_{p,p'}$ precedentes es «compatible con todo
cambio de base» $S' \longrightarrow S$ (es decir, corresponde a un
homomorfismo de secciones cartesianas de la categoría fibrada que se
considera sobre la categoría $\underline{T}_{/P\times P}$ de los pares
de flechas $p,p' : S \rightrightarrows P$ con destino $P$). El dato del
sistema de los $\varphi_{p,p'}$ que satisfacen esta condición equivale
también, por supuesto, al dato de $\varphi_{p,p'}$ en el caso
«universal», correspondiente al caso en que $S = P\times P$ y $p,p'$
son, respectivamente, las dos proyecciones $\mathrm{pr}_i$ ($i=1,2$)
de $P\times P$ en $P$. Se interpreta entonces como un isomorfismo de
bitorsors bajo $(G_{P\times P},G_{P\times P})$:

\begin{equation}\tag{1.1.3.2}
\varphi \ :\ \mathrm{pr}_1^*(E)\ \ \mathrm{pr}_2^*(E)\ \overset{\sim}{\longrightarrow}\ \pi^*(E) \qquad ,
\end{equation}
donde

\begin{equation*}
\pi \ :\ P\times P \ \longrightarrow\ P
\end{equation*}
es el morfismo de objetos de $\underline{T}$ que define la estructura
multiplicativa del grupo $P$. Sin embargo, las propiedades de
compatibilidad que siguen se enuncian más cómodamente en la forma de los
datos (1.1.3.1), evitando recurrir a diagramas menos próximos a la
intuición.


%%%%%%%%%%  scan idx 150  |  folio 139  %%%%%%%%%%

1.1.4.  Al escribir las relaciones (1.1.3.0), que nos permitieron
expresar la ley de composición dada sobre $E$ en términos de
isomorfismos de bitorsors (1.1.3.1), solo utilizamos la asociatividad de
esta ley de composición en el caso del producto de tres «puntos» de $E$
uno de los cuales, al menos, proviene de un punto de $G$. Si se quiere
expresar la asociatividad completa en términos de los
$\varphi_{p,p'}$, se llega a escribir, para tres puntos $p,p',p''$ de
$P$ (con valores en un objeto $S$ de $\underline{T}$), que el diagrama
siguiente es conmutativo:

\begin{equation}\tag{1.1.4.1}
\begin{tikzcd}[column sep=huge, row sep=large]
 & E_{pp'}E_{p''} \arrow[dr, "{\varphi_{pp',p''}}"] & \\
E_{p}E_{p'}E_{p''} \arrow[ur, "{\varphi_{p,p'}\wedge id}"] \arrow[dr, "{id\wedge\varphi_{p',p''}}"'] & & E_{pp'p''} \\
 & E_{p}E_{p'p''} \arrow[ur, "{\varphi_{p,p'p''}}"'] &
\end{tikzcd}
\end{equation}
% ESTRUCTURA: Rombo, 4 nodos, 4 flechas, todas flechas simples continuas con una sola punta rellena;
%   sin ganchos, sin cabezas dobles, sin tildes de isomorfismo ni trazos discontinuos. NODOS (4
%   exactamente): N1 (izquierda/vértice oeste, en la misma línea de base que la etiqueta impresa
%   «(1.1.4.1)» situada en el margen izquierdo): E_p E_p' E_p'' [E con subíndice p, E con subíndice p',
%   E con subíndice p'', yuxtapuestos]. N2 (arriba/vértice norte, por encima y a la derecha de N1):
%   E_pp' E_p''. N3 (abajo/vértice sur, por debajo y a la derecha de N1): E_p E_p'p''. N4
%   (derecha/vértice este): E_pp'p''. FLECHAS (4 exactamente): A1: N1 -> N2, flecha recta continua
%   ascendente hacia la derecha, etiqueta encima/a la izquierda del cuerpo: \varphi_{p,p'}\wedge id.
%   A2: N2 -> N4, flecha recta continua descendente hacia la derecha, etiqueta encima/a la derecha:
%   \varphi_{pp',p''}. A3: N1 -> N3, flecha recta continua descendente hacia la derecha, etiqueta
%   debajo/a la izquierda, ligeramente superpuesta a la punta en la fuente: id\wedge\varphi_{p',p''}.
%   A4: N3 -> N4, flecha recta continua ascendente hacia la derecha, etiqueta encima/a la derecha:
%   \varphi_{p,p'p''}. Las flechas entrantes A2 y A4 terminan en N4 desde la izquierda. Un punto final
%   suelto «.» aparece en el extremo derecho del área del diagrama, muy a la derecha y algo por debajo
%   de N4, al nivel de la fila N3.
En este diagrama hemos escrito el producto contraído de bitorsors sin
los signos $\overset{G_S}{\times}$, entendiéndose que la multiplicación
contraída de dos bitorsors $Q.R$ se efectúa utilizando (para pasar al
cociente en $Q\times R$) la operación del grupo estructural $G$
\underline{a derecha sobre} $Q$ \underline{y a izquierda sobre} $R$.
Se sabe, por otra parte, que esta operación de producto contraído es
asociativa salvo un isomorfismo canónico [2], y efectuamos la
identificación $(E_{p}E_{p'})E_{p''} \simeq
E_{p}(E_{p'}E_{p''})$, escribiendo simplemente ambos miembros como
$E_{p}E_{p'}E_{p''}$. En lo que sigue de esta exposición utilizaremos
convenciones evidentes de esta naturaleza sin señalarlas expresamente
cada vez.

1.1.4.2.  Del mismo modo que el dato de los isomorfismos
$\varphi_{p,p'}$ (1.1.3.1) se reduce al caso «universal»
$S=P\times P$, $p=\mathrm{pr}_1,p' = \mathrm{pr}_2$ considerado en (1.1.3.2), también
%%%%%%%%%%  scan idx 151  |  folio 140  %%%%%%%%%%
para verificar las condiciones de asociatividad (1.1.4.1) podemos
limitarnos al caso «universal» en que \sisi{$S=P\times P\times P$}{S=P$\times$P$\times$P},
\sisi{$p=\mathrm{pr}_1$}{$p=pr_1$},\sisi{$p'=\mathrm{pr}_2$}{$p'=pr_2$},\sisi{$p''=\mathrm{pr}_3$}{$p''=pr_3$}. Dejamos al lector el cuidado de escribir
el diagrama en cuestión con las notaciones que prefiera.
